\clearpage\SourceLeaf{247}\ArithmeticHead{162}{L}{L. SCHLESINGER, ÜBER GAUSS' FUNKTIONENTHEORETISCHE ARBEITEN.}
\noindent\hspace*{\parindent}Die \textit{Demonstratio tertia} beruht auf funktionentheoretischen Betrachtungen
über die Integration einer Funktion einer komplexen Veränderlichen. Es
werde für
\[
X=x^m+A_1x^{m-1}+\cdots+A_m,
\]
wo die \(A_1,\ldots,A_m\) reelle Grössen bedeuten \((A_m\ne0)\), \(x=re^{\varphi i}\) und \(X=t+ui\)
gesetzt. Ferner sei
\[
F(x)=x\frac{X'}{X}=P+Qi,\quad X'=\frac{dX}{dx}.
\]
\textsc{Gauss} betrachtet dann den Ausdruck
\[
\begin{aligned}
y&=\frac{\partial P}{\partial r}=\frac{\partial}{\partial r}
\left\{\frac{r\left(t\frac{\partial t}{\partial r}+u\frac{\partial u}{\partial r}\right)}{t^2+u^2}\right\}\\
&=\frac1r\frac{\partial Q}{\partial\varphi}=\frac1r\frac{\partial}{\partial\varphi}
\left\{\frac{r\left(t\frac{\partial u}{\partial r}-u\frac{\partial t}{\partial r}\right)}{t^2+u^2}\right\}
\end{aligned}
\]
und zeigt zunächst im Art. 1, dass für hinreichend grosse Werte von \(R\), insbesondere
für \(R\geqq\max\sqrt[\nu]{m|A_\nu|\sqrt2}\), die Ausdrücke
\[
r\left(t\frac{\partial t}{\partial r}+u\frac{\partial u}{\partial r}\right)\text{ und }t^2+u^2,
\]
also auch \(P\) für \(r=R\) wesentlich positiv sind, was ja übrigens aus dem
asymptotischen Verhalten von \(X\) und \(F(x)\)
\[
\lim_{x\to\infty}X\mathbin{.}x^{-m}=1,\quad\lim_{x\to\infty}(P+Qi)=m
\]
unmittelbar folgt. Nun sagt \textsc{Gauss} (\textit{Anzeige}, Werke III, S. 109): wenn im
Innern des Kreises mit dem Mittelpunkt \(0\) und dem Halbmesser \(R\), also für
\(0\leqq r<R,\ 0\leqq\varphi<2\pi\) kein Wertesystem \(r,\varphi\) vorhanden wäre, für das
gleichzeitig \(t=0,u=0\) wird, so folgt, dass \(t^2+u^2\) für jede Werte von
\(r\) und \(\varphi\) zwischen den angegebenen Grenzen immer positiv, und folglich \(y\)
immer endlich werden muss«\footnote{Ebenso im Art. 2 der Abhandlung selbst, Werke III, S. 61.}. Das will besagen, dass dann \(t^2+u^2\) oberhalb
einer angebbaren positiven Schranke und folglich \(y\) unterhalb einer ebensolchen
Schranke liegen muss. Die auf \(t^2+u^2\) bezügliche Behauptung bedarf
aber noch eines besondern Beweises, indem nämlich im Sinne des vierten Einwandes
gegen den Beweis von \textsc{D'Alembert}\footnote{In der \textit{Dissertation}, Werke III, S. 10, ad 4, vergl. oben S. 58, 59. Ein Hinweis auf diese »Lücke« in der \textit{Demonstr. tertia} findet sich bei A. \textsc{Fraenkel}, a. a. O. \textsuperscript{74)}, S. 13.} gezeigt werden muss, was \textsc{Gauss}
daselbst ausspricht, dass nämlich eine ganze rationale Funktion ihre obere
und untere Schranke in einem abgeschlossenen Bereich wirklich annimmt.\par
\clearpage\SourceLeaf{248}\ArithmeticHead{163}{R}{VIIIa. DER DRITTE WURZELEXISTENZBEWEIS.}
\noindent Setzt man das als richtig voraus, so verläuft der Beweis weiter wie folgt:
Wenn \(y\) beschränkt wäre, so hätte das über die abgeschlossene Kreisscheibe
\(C,r\leqq R\), erstreckte Doppelintegral
\begingroup\makeatletter\tagsleft@true\let\veqno\leqno\makeatother
\begin{equation*}
\iint\limits_C y\,dr\,d\varphi\tag{2}
\end{equation*}
\endgroup
\noindent einen wohlbestimmten endlichen Wert \(\Omega\), zu dessen Ermittlung die Integration
nach den beiden Veränderlichen \(r,\varphi\) in beliebiger Reihenfolge ausgeführt
werden dürfte. Nun ergibt sich, wenn man erst nach \(\varphi\) integriert, durch Betrachtung
des unbestimmten Integrals, dass für ein beliebiges \(r\)
\begingroup\makeatletter\tagsleft@true\let\veqno\leqno\makeatother
\begin{equation*}
\int_0^{2\pi}\frac{\partial P}{\partial r}\,d\varphi=0,\tag{3}
\end{equation*}
\endgroup
\noindent also auch \(\Omega=0\) ist, dagegen erhält man, wenn man erst nach \(r\) und dann
nach \(\varphi\) integriert
\[
\Omega=\int_0^{2\pi}(P(R,\varphi)-P(0,\varphi))\,d\varphi,
\]
es wäre also
\begingroup\makeatletter\tagsleft@true\let\veqno\leqno\makeatother
\begin{equation*}
\frac1{2\pi}\int_0^{2\pi}P(R,\varphi)\,d\varphi=P(0,\varphi),\tag{4}
\end{equation*}
\endgroup
\noindent was, da \(P(R,\varphi)\) wesentlich positiv, \(P(0,\varphi)\) aber gleich \(0\) ist, einen Widerspruch
involviert. Wir bemerken, dass die Gleichung (3) den sogenannten
\textsc{Gauss}schen Satz, die Gleichung (4) den sog. Mittelwertsatz der Potentialtheorie
darstellt, angewandt auf den Kreis \(r=R\) und die Potentialfunktion \(P\)\footnote{Vergl. M. \textsc{Bôcher}, Bulletin of the American Math. Soc. (2) 1, 1894, S. 205ff.}.
Wenn wir im Sinne der oben (S. 156) bei der \textit{Theoria attractionis} gemachten
Bemerkung annehmen dürfen, dass diese Sätze, die ja unmittelbare Folgerungen
aus dem Satze von der scheinbaren Grösse sind, \textsc{Gauss} schon 1813
bekannt waren, so würde man wohl vermuten können, dass er eben durch
diese Sätze auf die Betrachtung des Doppelintegrals (2) hingeleitet worden
ist. Aber es bedarf noch der Erklärung, weshalb \textsc{Gauss} gerade die Funktion
\(F(x)\) und nicht etwa\footnote{Siehe \textsc{Bôcher} a. a. O. und American Journal of Mathem. 17, 1895, S. 260ff.} die einfacher zu behandelnden \(x^mX^{-1}\) oder \(X^{-1}\) seiner
Betrachtung zugrunde legte. Diese Erklärung ergibt sich, wenn man unter\par
\HistoricalFoot{}{21\textsuperscript{*}}
\clearpage\SourceLeaf{249}\ArithmeticHead{164}{L}{L. SCHLESINGER, ÜBER GAUSS' FUNKTIONENTHEORETISCHE ARBEITEN.}
\noindent dem Doppelintegral (2) \(P\) durch die monogene Funktion \(P+Qi=F(x)\)
ersetzt, also \(\displaystyle\iint\frac{\partial F(x)}{\partial r}\,dr\,d\varphi\) betrachtet und die Integration, wie es \textsc{Gauss} im
Art. 4 der Abhandlung tut, über das durch die Ungleichungen \(k\leqq r\leqq l\),
\(\kappa\leqq\varphi\leqq\lambda\) gegebene Gebiet \(G\) erstreckt. Da \(\frac{\partial F}{\partial r}=-i\frac1r\frac{\partial F}{\partial\varphi}\) ist, hat man
dann
\[
\iint\limits_G\frac{\partial F}{\partial r}\,dr\,d\varphi
+i\iint\limits_G\frac1r\frac{\partial F}{\partial\varphi}\,dr\,d\varphi=0.
\]
Liegt nun in \(G\) keine Nullstelle von \(X\), so kann man in dem ersten Integral
die Integration nach \(r\), in dem zweiten die nach \(\varphi\) ausführen und erhält
\begingroup\makeatletter\tagsleft@true\let\veqno\leqno\makeatother
\begin{equation*}
\int_\kappa^\lambda[F(le^{\varphi i})-F(ke^{\varphi i})]\,d\varphi
+i\int_k^l\frac1r[F(re^{\lambda i})-F(re^{\kappa i})]\,dr=0.\tag{5}
\end{equation*}
\endgroup
\noindent Beachtet man, dass \(dx=e^{\varphi i}(dr+ir\,d\varphi)\) ist, und setzt
\[
x^{-1}F(x)=X^{-1}X'=\frac{d\log X}{dx}=L(x),
\]
so lässt sich die linke Seite der Gleichung (5) in der Form
\begingroup\makeatletter\tagsleft@true\let\veqno\leqno\makeatother
\begin{equation*}
\frac1i\left\{\begin{aligned}
&\int_\kappa^\lambda L(ke^{\varphi i})\,dx+\int_\lambda^\kappa L(le^{\varphi i})\,dx\\
&+\int_k^l L(re^{\kappa i})\,dx+\int_l^k L(re^{\lambda i})\,dx
\end{aligned}\right\}\tag{6}
\end{equation*}
\endgroup
\noindent oder kürzer
\[
\frac1i\int\limits_{\overline G}L(x)\,dx
\]
schreiben, wo \(\overline G\) die im positiven Sinne durchlaufene Begrenzung des Gebiets
\(G\) andeuten soll. Dieser Ausdruck stellt die Änderung dar, die \(\frac1i\log X\) erfährt,
wenn \(x\) die Begrenzung von \(G\) im positiven Sinne durchläuft. Wie
aus dem Briefe an \textsc{Bessel} vom 18. Dezember 1811 (Werke X, 1, S. 365, vergl.
oben S. 152) hervorgeht, wusste \textsc{Gauss}, dass der Logarithmus bei positiver Umkreisung
einer Nullstelle um \(2\pi i\) wächst, dass also das Integral (6) den Wert
\(2p\pi\) haben muss, wenn im Gebiete \(G\) genau \(p\) Nullstellen von \(X\) liegen.
Für \(k=0,l=R,\kappa=0,\lambda=2\pi\), wo also \(G\) in die Kreisperipherie \(C\) übergeht,
ist aber, da \(L(x)\) asymptotisch gleich \(mx^{-1}\) ist,
\[
\int\limits_C L(x)\,dx=2m\pi i,
\]
und somit \(m\) die Anzahl der einfach zu zählenden Nullstellen von \(X\) auf der\par
\clearpage\SourceLeaf{250}\ArithmeticHead{165}{R}{VIIIa. DER DRITTE WURZELEXISTENZBEWEIS.}
\noindent Kreisscheibe \(C\). Dass dieser Gedankengang in der Tat dem von \textsc{Gauss} befolgten
entspricht, geht einmal daraus hervor, dass er in der Fussnote zum
Art. 2 (Werke III, S. 62) den Wert des Integrals (6) für den Fall des Kreises
\(C\) tatsächlich gleich \(2m\pi\) angibt, andererseits aber auch aus einer Stelle des
Briefes an \textsc{Drobisch} (Werke X, 1, S. 107), wo er von dem dritten Beweise
sagt: »Auch ist derselbe im höchsten Grade sinnlich zu machen, was ich aber
dort für ein hors d'oeuvre hätte halten müssen, wo ich alles in rein algebraischer
Form darstellen wollte. Übrigens aber finde ich den \textsc{Cauchy}schen
Beweis\footnote{Gemeint ist der Beweis in der \textit{Analyse algébrique}, 1821, S. 329, vergl. dazu die Bemerkung Werke X, 1, S. 109, 110 und oben S. 59, 60.} in seiner Grundidee eben so elegant. \mbox{.\,.\,.\,.\,.} Aber am lehrreichsten
ist es wohl, beide Grundideen zu entwickeln, wo in der Tat die Betrachtung
der gleichsam geometrischen Bedeutung beider für den Verstand etwas recht
ergötzliches hat«. Diese »geometrische Bedeutung« des dritten Beweises dürfte
wohl im wesentlichen mit dem oben skizzierten Gedankengang übereinstimmen.
Übrigens gibt \textsc{Gauss} einen direkten Hinweis auf die Integration im komplexen
Gebiete im Art. 4, wenn er den Widerspruch beschreibt, in den man verstrickt
wird, wenn man ein Integral über eine unendlich werdende Funktion,
wie z. B. \(x^{-2}\), nach den gewöhnlichen Regeln entwickeln will, und hinzufügt:
»was aber diese und ähnliche analytische Paradoxien bedeuten, das soll bei
einer anderen Gelegenheit ausführlich verfolgt werden«; freilich ist dieser Hinweis
nur für denjenigen verständlich, der bereits weiss, dass man der Unendlichkeitsstelle
durch geeignete Wahl des Integrationsweges in der komplexen
Ebene ausweichen kann.
\par
Da die Änderung von \(\frac1i\log X\) beim Durchlaufen der Kreisperipherie \(C\)
nichts anderes ist, als die Änderung von \(\operatorname{arctg}\frac ut\) oder also die von \(\arg X\), so
berührt sich der dritte Beweis mit der am Schluss des Abschnitts III (oben
S. 61, 62) angedeuteten Auffassung des Beweises der \textit{Dissertation}. Auf einen
solchen Zusammenhang hat \textsc{Gauss} selbst in einer Vorlesung \textit{Über die Theorie
der imaginären Grössen} hingewiesen, die er von Neujahr bis Ostern 1840 gehalten
hat, und von der eine Ausarbeitung sich im Gaussarchiv befindet\footnote{Vergl. hierzu die in der Fussnote \textsuperscript{80)} zitierte Arbeit von A. \textsc{Fraenkel}. Darstellungen des dritten Beweises findet man in der in der Fussnote \textsuperscript{239)} genannten Arbeit von \textsc{Bôcher} sowie in E. \textsc{Goursats} \textit{Cours d'Analyse} I, 1910, S. 240. Zur Geschichte der Integration durch imaginäres Gebiet und \textsc{Gauss'} Einfluss auf die Entwicklung dieser Theorie vergleiche man P. \textsc{Stäckels} Aufsätze, Bibliotheca Math. (3) 1, 1900, S. 109; 2, 1901, S. 111, sowie Bemerkungen desselben zu seiner Ausgabe von \textsc{Cauchys} Abhandlung von 1825, \textsc{Ostwalds} Klassiker Nr. 112, S. 67—72.}.\par
